% GENERATED FILE. Edit canonical lessons, then run npm run generate:books.
% canonical-source-hash: fnv1a64:e6ea28d24731c4a8
% canonical-lessons: ZH-C311-diaocha, ZH-C311-jingche, ZH-W311-jing-alert, ZH-C311-zhui, ZH-C311-jingbao, ZH-C311-jiaojing

\chapter{The Police Car}
\label{ch:zh-the-police-car}
\hlchaptermodality{311}

\begin{chapteropening}
\textbf{By the end of this chapter:} \emph{I can talk about an investigation, a police car, a chase, an alarm and the traffic police.}

Complete the last lesson of chapter 311: I can talk about an investigation, a police car, a chase, an alarm and the traffic police.
\end{chapteropening}

\section[diàochá]{diàochá — to investigate}

\label{lesson:ZH-C311-diaocha}



\begin{quote}
Before the new one: say the Mandarin for to tell a lie, then the Mandarin for to own up to something.
\end{quote}

\subsection*{You'll want to know: diàochá}
\emph{diàochá} — \textquotedblleft{}to investigate\textquotedblright{}.

\begin{quote}
\emph{jǐngchá zhèngzài diàochá zhège shìgù} — The police are investigating the accident.
\end{quote}

Its characters come later; for now it is a word for the ear.

\begin{cousinweb}
\emph{diàochá} is built from \emph{diào} 'to choose' and \emph{chá}, both inherited from Middle Chinese, \emph{dewH} and \emph{dzrae} (Guangyun, in Baxter's transcription), which glosses the first 'to select' and the second 'wood floating in water'; how the pair came to mean 'to investigate' is not settled.
\end{cousinweb}

\begin{grammarlens}[title={What you've built}]
One more word for telling what happened.
\end{grammarlens}

\subsection*{Your turn}
\begin{itemize}
\raggedright
  \item \emph{Say it:} \emph{diàochá}
  \item \emph{Say it:} \emph{diàochá}, once more
  \item \emph{Say it:} \emph{jǐngchá zhèngzài diàochá zhège shìgù} — The police are investigating the accident.
  \item \emph{Recall:} say the Mandarin for to tell a lie, then the Mandarin for to own up to something, then say \emph{diàochá} again
\end{itemize}

\begin{tcolorbox}[breakable,colback=teal!4,colframe=teal!35!black,title={Before you move on}]
What does diàochá mean? (\textquotedblleft{}To investigate\textquotedblright{}.) Say it once more.
\end{tcolorbox}
\section[\texorpdfstring{\zh{警车} (jǐngchē)}{警车 (jǐngchē)}]{\zh{警车} (jǐngchē) — police car}

\label{lesson:ZH-C311-jingche}



\begin{quote}
Before the new one: say the Mandarin for to own up to something, then the Mandarin for to investigate a case.
\end{quote}

\subsection*{You'll want to know: \zh{警车}}
\textbf{\zh{警车}} — \emph{jǐngchē} — \textquotedblleft{}police car\textquotedblright{}.

\begin{quote}
\textbf{\zh{小偷看到警车就走了}} \emph{xiǎotōu kàndào jǐngchē, jiù zǒu le} — The thief saw the police car and left.
\end{quote}

The new character, in pieces: \zh{警} — nineteen strokes in three pieces: top left, the grass piece of \zh{花}, then a two-stroke wrap with \zh{口} inside; top right, the striking piece of \zh{救}; across the foot a dot, three bars and \zh{口}. The next lesson puts it in your hand.

\begin{cousinweb}
\textbf{\zh{警车}} (\emph{jǐngchē}) is built from \zh{警} 'to warn' and \zh{车} 'a vehicle', both inherited from Middle Chinese, \emph{kjaengX} and \emph{tsyhae} (Guangyun, in Baxter's transcription), which glosses the first 'to wake' and 'to guard against'.
\end{cousinweb}

\begin{grammarlens}[title={What you've built}]
One more word for telling what happened.
\end{grammarlens}

\subsection*{Your turn}
\begin{itemize}
\raggedright
  \item \emph{Say it:} \emph{jǐngchē}
  \item \emph{Say it:} \emph{jǐngchē}, once more
  \item \emph{Say it:} \emph{xiǎotōu kàndào jǐngchē, jiù zǒu le} — The thief saw the police car and left.
  \item \emph{Recall:} say the Mandarin for to own up to something, then the Mandarin for to investigate a case, then say \emph{jǐngchē} again
\end{itemize}

\begin{tcolorbox}[breakable,colback=teal!4,colframe=teal!35!black,title={Before you move on}]
What does \zh{警车} mean? (\textquotedblleft{}Police car\textquotedblright{}.) Say it once more.
\end{tcolorbox}
\section[\texorpdfstring{\zh{警} (jǐng)}{警 (jǐng)}]{\zh{警} — three pieces: top left, top right and the foot}

\label{lesson:ZH-W311-jing-alert}



\begin{quote}
The word first, then the character inside it.

\begin{itemize}
\raggedright
  \item \emph{Recall:} say \emph{jǐngchē} — \textbf{R1}, one lesson back, before it decays
\end{itemize}

\begin{quote}
\zh{花} \zh{救}
\end{quote}

\emph{Read it:} each one aloud, with nothing said first — \emph{huā}, \emph{jiù}
\end{quote}

\subsection*{Script — three pieces: top left, top right and the foot}
\hlblockfigure{\detokenize{figures/ZH-W311-jing-alert-filmstrip.pdf}}{How \zh{警} is written, stroke by stroke}

\begin{quote}
\zh{警}
\end{quote}

\textbf{Three pieces: two side by side on top, one across the foot.} Count the pieces, not the strokes.

\textbf{Top left, eight strokes.} The grass piece of \zh{花}: a bar and two short verticals. Under it a short sweep to the left and a stroke that runs right, turns down and hooks back. Inside, \zh{口}.

\textbf{Top right, four strokes.} The striking piece of \zh{救}: a sweep to the left, a bar, a long sweep to the lower left and a long stroke falling to the right.

\textbf{The foot, seven strokes.} A dot. A long bar across, then two shorter bars. Then \zh{口} under them.

8 + 4 + 7 = nineteen.

\subsection*{Your turn}
\begin{itemize}
\raggedright
  \item \emph{Write it:} \zh{警}, all nineteen, with the model in front of you
  \item \emph{Cover:} write it again from memory, and count the pieces, not the strokes
\end{itemize}

\begin{tcolorbox}[breakable,colback=teal!4,colframe=teal!35!black,title={Before you move on}]
How many pieces in \zh{警}? (\textbf{Three}.) Which piece sits across the foot? (\textbf{A dot, three bars and \zh{口}}.) Where have you made the top right piece? (\textbf{In \zh{救} and \zh{故}}.)

Source: \href{https://raw.githubusercontent.com/chanind/hanzi-writer-data/68d10a4b21150cae5e1ebbd223eed289cf32d90c/data/\%E8\%AD\%A6.json}{Hanzi Writer Data \zh{警}, pinned stroke paths and medians}.
\end{tcolorbox}
\section[zhuī]{zhuī — to chase; to run after}

\label{lesson:ZH-C311-zhui}



\begin{quote}
Before the new one: say the Mandarin for to investigate a case, then the Mandarin for a police car.
\end{quote}

\subsection*{You'll want to know: zhuī}
\emph{zhuī} — \textquotedblleft{}to chase; to run after\textquotedblright{}.

\begin{quote}
\emph{jǐngchá zài zhuī xiǎotōu} — The police are chasing the thief.
\end{quote}

Its characters come later; for now it is a word for the ear.

\begin{cousinweb}
\emph{zhuī} is inherited from Middle Chinese \emph{trwij} (Guangyun, in Baxter's transcription), which glosses it 'to chase' and 'to follow'.
\end{cousinweb}

\begin{grammarlens}[title={What you've built}]
One more word for telling what happened.
\end{grammarlens}

\subsection*{Your turn}
\begin{itemize}
\raggedright
  \item \emph{Say it:} \emph{zhuī}
  \item \emph{Say it:} \emph{zhuī}, once more
  \item \emph{Say it:} \emph{jǐngchá zài zhuī xiǎotōu} — The police are chasing the thief.
  \item \emph{Recall:} say the Mandarin for to investigate a case, then the Mandarin for a police car, then say \emph{zhuī} again
\end{itemize}

\begin{tcolorbox}[breakable,colback=teal!4,colframe=teal!35!black,title={Before you move on}]
What does zhuī mean? (\textquotedblleft{}To chase; to run after\textquotedblright{}.) Say it once more.
\end{tcolorbox}
\section[\texorpdfstring{\zh{警报} (jǐngbào)}{警报 (jǐngbào)}]{\zh{警报} (jǐngbào) — alarm; a warning siren}

\label{lesson:ZH-C311-jingbao}



\begin{quote}
Before the new one: say the Mandarin for a police car, then the Mandarin for to run after someone.
\end{quote}

\subsection*{You'll want to know: \zh{警报}}
\textbf{\zh{警报}} — \emph{jǐngbào} — \textquotedblleft{}alarm; a warning siren\textquotedblright{}.

\begin{quote}
\textbf{\zh{有警报大家快走}} \emph{yǒu jǐngbào, dàjiā kuài zǒu} — There's an alarm, everyone out, quick!
\end{quote}

\begin{cousinweb}
\textbf{\zh{警报}} (\emph{jǐngbào}) is built from \zh{警} 'to warn' and \zh{报} 'to report', both inherited from Middle Chinese, \emph{kjaengX} and \emph{pawH} (Guangyun, in Baxter's transcription): a warning sent out.
\end{cousinweb}

\begin{grammarlens}[title={What you've built}]
One more word for telling what happened.
\end{grammarlens}

\subsection*{Your turn}
\begin{itemize}
\raggedright
  \item \emph{Say it:} \emph{jǐngbào}
  \item \emph{Say it:} \emph{jǐngbào}, once more
  \item \emph{Say it:} \emph{yǒu jǐngbào, dàjiā kuài zǒu} — There's an alarm, everyone out, quick!
  \item \emph{Recall:} say the Mandarin for a police car, then the Mandarin for to run after someone, then say \emph{jǐngbào} again
\end{itemize}

\begin{tcolorbox}[breakable,colback=teal!4,colframe=teal!35!black,title={Before you move on}]
What does \zh{警报} mean? (\textquotedblleft{}Alarm; a warning siren\textquotedblright{}.) Say it once more.
\end{tcolorbox}
\section[jiāojǐng]{jiāojǐng — traffic police officer}

\label{lesson:ZH-C311-jiaojing}



\begin{quote}
Before the new one: say the Mandarin for to run after someone, then the Mandarin for a warning siren.
\end{quote}

\subsection*{You'll want to know: jiāojǐng}
\emph{jiāojǐng} — \textquotedblleft{}traffic police officer\textquotedblright{}.

\begin{quote}
\emph{jiāojǐng shuō wǒ kāichē tài kuài} — The traffic officer said I was driving too fast.
\end{quote}

Its characters come later; for now it is a word for the ear.

\begin{cousinweb}
\emph{jiāojǐng} is built from \emph{jiāo} 'to cross' and \zh{警} 'to warn', a short form of a longer name for the police of the roads; both are inherited from Middle Chinese, \emph{kaew} and \emph{kjaengX} (Guangyun, in Baxter's transcription), which glosses the first 'to join'.
\end{cousinweb}

\begin{grammarlens}[title={What you've built}]
One more word for telling what happened.
\end{grammarlens}

\subsection*{Your turn}
\begin{itemize}
\raggedright
  \item \emph{Say it:} \emph{jiāojǐng}
  \item \emph{Say it:} \emph{jiāojǐng}, once more
  \item \emph{Say it:} \emph{jiāojǐng shuō wǒ kāichē tài kuài} — The traffic officer said I was driving too fast.
  \item \emph{Recall:} say the Mandarin for to run after someone, then the Mandarin for a warning siren, then say \emph{jiāojǐng} again
\end{itemize}

\begin{tcolorbox}[breakable,colback=teal!4,colframe=teal!35!black,title={Before you move on}]
What does jiāojǐng mean? (\textquotedblleft{}Traffic police officer\textquotedblright{}.) Say it once more.
\end{tcolorbox}
